% =====================================================================
\section*{Введение}
\phantomsection
\addcontentsline{toc}{section}{Введение}

В отчете представлено описание выполнения лабораторной работы №1 по дисциплине «Теория алгоритмов».

Необходимо реализовать двумерный клеточный автомат с окрестностью фон Неймана в соответствии с полученным номером №. Граничные условия --- торроидальные. Пользователь определяет ширину поля и количество итераций. Необходимо учесть возможность ввода различных начальных условий (как вручную, так и случайным образом) по выбору пользователя. Реализация возможна в консоли или в графике. В итоге необходимо проанализировать свой клеточный автомат в отчете (его поведение, паттерны, «сходимость», периодичность, классификация по Вольфраму и т.д.).
№ = номер варианта $\cdot$ 11 $\cdot$ год рождения $\cdot$ месяц $\cdot$ день.

Для варианта 43 номер правила равен $\text{№} = 43 \cdot 11 \cdot 2006 \cdot 11 \cdot 2 = 20\,874\,436.$

Программа реализована на языке Java в виде консольного приложения. % В разделе~\ref{sec:math} приведено математическое описание автомата и правила №, в разделе~\ref{sec:impl} --- особенности реализации, в разделе~\ref{sec:analysis} --- анализ поведения автомата, в разделе~\ref{sec:results} --- примеры работы программы.
